\section{خلاصة الفصل}
تناول هذا الفصل تصميم بنية نظام التفاعل والتحكم الصوتي للذراع الروبوتية ذات الخمس درجات حرية، ليمثل الطبقة العليا الإشرافية في الهيكلية الهرمية للمشروع. تم البدء بتحديد البنية العامة لتدفق الإشارة الصوتية بدءاً من الالتقاط الفيزيائي عبر وحدة التقاط الصوت وحتى المعالجة الرقمية واستخلاص النص المكتوب.

وقد استعرض الفصل آلية عمل خوارزمية \textenglish{Vosk} كخيار هندسي كفء يتيح التعرف الذاتي على الكلام دون اتصال بالشبكة وبأقل استهلاك للموارد العتادية للأنظمة المدمجة، مع توضيح التكامل الرياضي بين النموذج الصوتي والنموذج اللغوي لتوليد الكلمات بدقة. كما تم تفصيل منهجية ربط النصوص المستخرجة مباشرة بالأفعال الحركية الكارتزية وشرح تتابعها عبر مخطط تدفق مبسط، مما يضمن تحديث شعاع الموضع المرجعي وتمريره اللحظي إلى طبقات التحكم عالي المستوى (\textenglish{HLC}) ومنخفض المستوى (\textenglish{LLC}) للقيادة الآمنة والمستقرة لمفاصل الذراع، وبذلك تصبح المنظومة متكاملة وجاهزة لمرحلة التطبيق والتحقق العملي.